\chapter{A Chabauty--Coleman bound for surfaces}

% archon:covers MR4665779AChabautyColemanBoundForSurfaces.lean MR4665779AChabautyColemanBoundForSurfaces/Basic.lean

This chapter is the complete informal dependency roadmap for Caro--Pasten's
surface analogue of Coleman's bound.  It deliberately gives placeholder Lean
names for the mathematical declarations that are not yet scaffolded; the only
current Lean declaration is the bootstrap constant \(\lean{hello}\).

\section{Introduction}

\subsection{Bootstrap coverage}

\begin{definition}[Bootstrap constant]
  \label{def:bootstrap_hello}
  \lean{hello}
  The current Lean scaffold contains the constant \(\mathrm{hello}\), whose
  value is the string ``world''.  It carries no mathematical content and is
  included only to keep the Lean--blueprint correspondence exact.
\end{definition}

\subsection{Results}

\begin{theorem}[Complete Caro--Pasten theorem package]
  \label{thm:cp_project_goal}
  \lean{MR4665779.ProjectGoal}
  \uses{def:bootstrap_hello, thm:cp_main_local, thm:cp_main_qp, thm:cp_number_field_bound, thm:cp_intro_surface_bounds, cor:cp_dim3_density_application, cor:cp_quadratic_points, cor:cp_quadratic_points_genus_three}
  % SOURCE: Caro--Pasten, Abstract and Sections 1, 9, 10 (read from references/caro-pasten-chabauty-coleman-surfaces-source/MR4665779 - A Chabauty-Coleman bound for surfaces.tex)
  % SOURCE QUOTE: We prove a result analogous to Coleman's bound in the case of a hyperbolic surface $X$ over a number field, embedded in an abelian variety $A$ of rank at most one.
  \textit{Source: Caro--Pasten, Abstract and Sections 1, 9, 10.}
  The project proves the complete Caro--Pasten blueprint package: the local
  Chabauty--Coleman bound for surfaces over finite extensions of
  \(\mathbb Q_p\), its \(\mathbb Q_p\) specialization, the number-field
  rational-point consequences, density-one applicability for abelian
  threefolds with endomorphism ring \(\mathbb Z\), and the symmetric-square
  applications bounding quadratic points on nonhyperelliptic curves.
\end{theorem}
\begin{proof}
  The theorem is the conjunction of the local theorem, its specialization and
  number-field transfer, the two rational-surface applications, Chavdarov's
  density corollary, and the two symmetric-square corollaries.  The bootstrap
  Lean constant is included only so the current Lean scaffold lies in the same
  dependency cone; it contributes no mathematical premise.
\end{proof}

\section{Notation}

\subsection{Global notation}

\begin{definition}[Geometric and local-field notation]
  \label{def:cp_notation}
  \lean{MR4665779.Notation}
  % SOURCE: Caro--Pasten, Section 2 (Notation), subsections \label{NotationFields}, \label{notationGeo}, \label{NotationLoc} (read from references/caro-pasten-chabauty-coleman-surfaces-source/MR4665779 - A Chabauty-Coleman bound for surfaces.tex)
  % SOURCE QUOTE: A variety over $L$ is a reduced $L$-scheme which is separated of finte type over $L$.
  \textit{Source: Caro--Pasten, Section 2 (NotationFields, notationGeo, NotationLoc).}
  A variety over a field \(L\) is a reduced separated \(L\)-scheme of finite
  type.  Curves and surfaces have pure dimensions \(1\) and \(2\).  For a smooth
  projective surface \(S\), \(c_1^2(S)\) is the self-intersection of a canonical
  divisor.  For a finite extension \(K/\mathbb Q_p\), write \(R\) for its ring
  of integers, \(\mathfrak p\) for its maximal ideal, \(\varpi\) for a
  uniformizer, \(\mathbf k\) for the residue field of cardinality \(q\),
  \(\mathfrak e\) for the ramification index, and normalize the absolute value
  by \(|\varpi|=q^{-1}\).
\end{definition}

\section{Curves and surfaces}

\subsection{Geometric interfaces}

\begin{definition}[Basic surface and abelian-variety interfaces]
  \label{def:cp_geometry_interfaces}
  \lean{MR4665779.GeometryInterfaces}
  \uses{def:cp_notation}
  % SOURCE: Caro--Pasten, Sections 3 and 5 (read from references/caro-pasten-chabauty-coleman-surfaces-source/MR4665779 - A Chabauty-Coleman bound for surfaces.tex)
  % SOURCE QUOTE: Given a geometrically irreducible, smooth, projective surface $S$ over $L$, we have a $\Z$-valued intersection product on divisors of $S$.
  \textit{Source: Caro--Pasten, Sections 3 and 5.}
  The formalization uses the standard interfaces for normalizations of curves,
  arithmetic and geometric genus, multiplicity and singularity order at curve
  points, divisor intersection on smooth projective surfaces, canonical divisors
  and nef/ample divisors, global Kähler differentials, abelian varieties,
  translations, invariant differentials, Tate-module characteristic
  polynomials, and specialization of cycles under good reduction.
\end{definition}

\subsection{Curves, surfaces, and point counting}

\begin{lemma}[Curve singularities and normalization]
  \label{lem:cp_curve_singularity_package}
  \lean{MR4665779.CurveSingularityPackage}
  \uses{def:cp_geometry_interfaces}
  % SOURCE: Caro--Pasten, Lemmas 3.1--3.5 (read from references/caro-pasten-chabauty-coleman-surfaces-source/MR4665779 - A Chabauty-Coleman bound for surfaces.tex)
  % SOURCE QUOTE: If $C$ is a projective irreducible curve over $k$, then $\gfrak_a(C)=\gfrak_g(C) + \sum_{x\in C(k)} \delta(C,x)$.
  \textit{Source: Caro--Pasten, Lemmas 3.1--3.5.}
  For an irreducible curve \(C\) over an algebraically closed field and
  \(x\in C(k)\), the singularity order \(\delta(C,x)\) is nonnegative and
  vanishes exactly at regular points.  If \(C\) is projective, then
  \[
    g_a(C)=g_g(C)+\sum_{x\in C(k)}\delta(C,x).
  \]
  If \(C\) lies on a regular surface at \(x\), the multiplicity is the largest
  \(n\) for which a local equation lies in \(\mathfrak m_{S,x}^n\), and
  \[
    \delta(C,x)\ge \frac12\,\operatorname{mult}_x(C)
      (\operatorname{mult}_x(C)-1).
  \]
\end{lemma}
\begin{proof}
  Nonnegativity and the regularity criterion follow from the normalization
  quotient.  Hironaka's genus formula identifies the total discrepancy between
  arithmetic and geometric genus with the sum of local \(\delta\)-invariants.
  The multiplicity description is the standard Hilbert--Samuel formula for a
  plane curve germ, and Hironaka's neighboring-points inequality gives the
  displayed lower bound after discarding nonnegative higher contributions.
\end{proof}

\begin{lemma}[Canonical genus bound]
  \label{lem:cp_canonical_genus_bound}
  \lean{MR4665779.CanonicalGenusBound}
  \uses{def:cp_geometry_interfaces, lem:cp_curve_singularity_package}
  % SOURCE: Caro--Pasten, Lemma 3.6 (read from references/caro-pasten-chabauty-coleman-surfaces-source/MR4665779 - A Chabauty-Coleman bound for surfaces.tex)
  % SOURCE QUOTE: Then we have $$\sum_{j=1}^\ell a_j\cdot (\gfrak_a(C_j)-1)\le c_1^2(S).$$
  \textit{Source: Caro--Pasten, Lemma 3.6.}
  Let \(S\) be a smooth projective irreducible surface, and let
  \(D=\sum_{j=1}^{\ell}a_jC_j\) be a positive nef canonical divisor.  Then
  \[
    \sum_{j=1}^{\ell}a_j\bigl(g_a(C_j)-1\bigr)\le c_1^2(S).
  \]
\end{lemma}
\begin{proof}
  Reorder the \(C_j\) so that the positive self-intersection components come
  first.  Adjunction gives
  \(2\sum a_j(g_a(C_j)-1)=D^2+\sum a_j C_j^2\).  For each positive
  self-intersection component, \((D-a_jC_j)\cdot C_j\ge0\), and nefness gives
  \(D\cdot C_j\ge0\); hence \(D^2\) bounds the positive contribution
  \(\sum a_j^2 C_j^2\).  Combining the two inequalities gives the result.
\end{proof}

\begin{lemma}[Branches and normalization fibres]
  \label{lem:cp_branches_package}
  \lean{MR4665779.BranchesPackage}
  \uses{lem:cp_curve_singularity_package}
  % SOURCE: Caro--Pasten, Lemmas 3.7--3.10 (read from references/caro-pasten-chabauty-coleman-surfaces-source/MR4665779 - A Chabauty-Coleman bound for surfaces.tex)
  % SOURCE QUOTE: there is a bijection between the set of branches of $C$ at $x$ and the points $y\in \nu^{-1}(x)$.
  \textit{Source: Caro--Pasten, Lemmas 3.7--3.10.}
  For a curve \(C\) through a smooth point \(x\) of a surface \(S\), the
  irreducible factors of a completed local equation are pairwise
  non-associated and define the branches at \(x\).  These branches are exactly
  the kernels of the completed maps associated to points of the normalization
  fibre \(\nu^{-1}(x)\), and
  \(\#\nu^{-1}(x)\le\delta(C,x)+1\).  If
  \(f(x_1,x_2)\in k[[x_1,x_2]]\) is irreducible and
  \(f(\phi_1(t),\phi_2(t))=0\), then
  \(\operatorname{ord}_t\phi_1\ge\operatorname{ord}_{x_2}f(0,x_2)\) and
  \(\operatorname{ord}_t\phi_2\ge\operatorname{ord}_{x_1}f(x_1,0)\).
\end{lemma}
\begin{proof}
  Analytic unramifiedness of the curve germ excludes associated distinct
  factors in the completed local ring.  The standard branch--normalization
  correspondence identifies branches with completed kernels at normalization
  preimages.  The number of branches is bounded by multiplicity, and
  multiplicity is bounded in terms of \(\delta\) by the singularity package.
  The order estimates are Płoski's theorem for irreducible plane power series.
\end{proof}

\begin{lemma}[General type and ampleness criteria]
  \label{lem:cp_general_type_package}
  \lean{MR4665779.GeneralTypePackage}
  \uses{def:cp_geometry_interfaces, lem:cp_curve_singularity_package}
  % SOURCE: Caro--Pasten, Lemmas 3.11--3.14 (read from references/caro-pasten-chabauty-coleman-surfaces-source/MR4665779 - A Chabauty-Coleman bound for surfaces.tex)
  % SOURCE QUOTE: If $S(\C)$ is hyperbolic, then $S$ is of general type.
  \textit{Source: Caro--Pasten, Lemmas 3.11--3.14.}
  A minimal smooth projective surface of general type has \(c_1^2\ge1\) and
  nef canonical sheaf.  If such a surface has no smooth rational curve of
  self-intersection \(-1\) or \(-2\), its canonical sheaf is ample.  A complex
  hyperbolic smooth projective surface is of general type.  If
  \(X\subset A\) is a smooth projective surface in a complex abelian variety,
  \(X\) is not an abelian surface, and \(X\) contains no elliptic curves, then
  \(X\) is of general type.
\end{lemma}
\begin{proof}
  The nefness and positivity are the standard minimal-surface facts.  For
  ampleness, Nakai--Moishezon reduces to excluding curves with zero
  canonical degree; the Hodge index theorem and adjunction turn such a curve
  into a forbidden rational \((-2)\)-curve.  Hyperbolicity implies general type
  for surfaces by the cited surface classification theorem.  Green--Yau
  hyperbolicity for subvarieties of abelian varieties without abelian-surface
  or elliptic-curve exceptions gives the final assertion.
\end{proof}

\begin{lemma}[Modified Weil bound for singular curves]
  \label{lem:cp_weil_bound_singular_curve}
  \lean{MR4665779.WeilBoundSingularCurve}
  \uses{lem:cp_curve_singularity_package}
  % SOURCE: Caro--Pasten, Lemmas 3.15--3.17 (read from references/caro-pasten-chabauty-coleman-surfaces-source/MR4665779 - A Chabauty-Coleman bound for surfaces.tex)
  % SOURCE QUOTE: $$\sum_{x\in D(\kk)} \sum_{j=1}^r \#\{y\in \widetilde{C}_j(\kk^{alg}) : \nu_j(y)=x\} \le (q+1)r + 2q^{1/2}\cdot \sum_{j=1}^r \gfrak_g(C_j).$$
  \textit{Source: Caro--Pasten, Lemmas 3.15--3.17.}
  Let \(D\) be a projective curve over a finite field \(\mathbf k\) of
  cardinality \(q\), and let \(C_1,\ldots,C_r\) be its geometric irreducible
  components with normalizations \(\widetilde C_j\).  Then
  \[
    \sum_{x\in D(\mathbf k)}\sum_{j=1}^r
      \#\{y\in\widetilde C_j(k^{alg})\mid \nu_j(y)=x\}
    \le (q+1)r+2q^{1/2}\sum_{j=1}^r g_g(C_j).
  \]
\end{lemma}
\begin{proof}
  After a finite field extension over which all components, singularities, and
  normalization preimages are rational, the excess
  \(A_D(L)-\#D(L)\) stabilizes.  The modified zeta function
  \(Z^*_{D,\mathbf k}\) satisfies a product identity over roots of unity with
  the zeta functions of the smooth normalizations.  Weil's bounds for smooth
  projective curves applied to those factors, followed by logarithmic
  differentiation and coefficient comparison at \(n=1\), gives the inequality.
\end{proof}

\section{\texorpdfstring{\(\omega\)}{omega}-integrality}

\begin{definition}[\(\omega\)-integral morphisms]
  \label{def:cp_omega_integral}
  \lean{MR4665779.OmegaIntegral}
  \uses{def:cp_geometry_interfaces}
  % SOURCE: Caro--Pasten, Section 4 (\label{Secwint}, subsection ``$\omega$-integral morphisms''; the definition is unnumbered prose) (read from references/caro-pasten-chabauty-coleman-surfaces-source/MR4665779 - A Chabauty-Coleman bound for surfaces.tex)
  % SOURCE QUOTE: Given $T$-schemes $X$ and $Y$ and a differential $\omega\in H^0(Y, \Omega^1_{Y/T})$, we say that a morphism of $T$-schemes $\phi:X\to Y$ is \emph{$\omega$-integral} if $\phi^\bullet(\omega)=0$.
  \textit{Source: Caro--Pasten, Section 4 (Secwint), ``$\omega$-integral morphisms''.}
  For a morphism of \(T\)-schemes \(\phi:X\to Y\), the pullback of
  differentials defines
  \[
    \phi^\bullet:H^0(Y,\Omega^1_{Y/T})\to H^0(X,\Omega^1_{X/T}).
  \]
  The morphism is \(\omega\)-integral when \(\phi^\bullet(\omega)=0\).
\end{definition}

\begin{lemma}[Functoriality and locality of \(\omega\)-integrality]
  \label{lem:cp_omega_integral_functorial_local}
  \lean{MR4665779.OmegaIntegralFunctorialLocal}
  \uses{def:cp_omega_integral}
  % SOURCE: Caro--Pasten, Lemmas 4.1 (LemmaCompositionInt) and 4.3 (LemmaCommLoc) (read from references/caro-pasten-chabauty-coleman-surfaces-source/MR4665779 - A Chabauty-Coleman bound for surfaces.tex)
  % SOURCE QUOTE: Then $(\psi\circ \phi)^\bullet = \phi^\bullet \circ \psi^\bullet$.
  \textit{Source: Caro--Pasten, Lemmas 4.1 and 4.3.}
  Pullback of differentials is strictly functorial:
  \((\psi\circ\phi)^\bullet=\phi^\bullet\circ\psi^\bullet\).  It is also
  compatible with local rings: the localization of \(\phi^\bullet(\omega)\) at
  \(x\) equals the pullback of the localized germ of \(\omega\) by
  \(\phi_x^\#\).
\end{lemma}
\begin{proof}
  The functorial statement follows by comparing the canonical maps
  \(\mathcal F\to f_*f^*\mathcal F\) and the maps from pulled-back
  differentials to relative differentials in the three commutative diagrams of
  Caro--Pasten.  The local statement is the same construction after localizing
  the sheaves of differentials and the morphism of rings at the chosen points.
\end{proof}

\begin{lemma}[Local jet expressions for integral differentials]
  \label{lem:cp_local_integral_expressions}
  \lean{MR4665779.LocalIntegralExpressions}
  \uses{def:cp_omega_integral, lem:cp_omega_integral_functorial_local}
  % SOURCE: Caro--Pasten, Lemmas 4.8--4.10 (LemmaGenOver, LemmaLocDiff, LemmaLocD) and Corollary 4.11 (Corom0) (read from references/caro-pasten-chabauty-coleman-surfaces-source/MR4665779 - A Chabauty-Coleman bound for surfaces.tex)
  % SOURCE QUOTE: There are local parameters $s_1,s_2\in \mfrak_{S,x}$ at $x$ such that $\phi^\#_\zeta(s_1)=z \bmod z^{m+1}$ and $\phi^\#_\zeta(s_2)=0$.
  \textit{Source: Caro--Pasten, Lemmas 4.8--4.10 and Corollary 4.11.}
  Let \(\phi:\operatorname{Spec}k[z]/(z^{m+1})\to S\) be a closed immersion
  supported at a smooth surface point \(x\).  There are local parameters
  \(s_1,s_2\) with \(\phi^\#(s_1)=z\), \(\phi^\#(s_2)=0\), and
  \(\ker\phi^\#=(s_1^{m+1},s_2)\).  If \(\phi\) is \(\omega\)-integral, then
  locally
  \[
    \omega=(G_1s_1^m+G_2s_2)\,ds_1+G_3\,ds_2.
  \]
  Consequently, if \(\omega_1\wedge\omega_2=F\,ds_1\wedge ds_2\), then
  \(F\in(s_1^m,s_2)\), and if \(x\) is not in the divisor of
  \(\omega_1\wedge\omega_2\), then \(m=0\).
\end{lemma}
\begin{proof}
  Surjectivity of the local map of a closed immersion lets one choose
  \(s_1\) mapping to \(z\); a second parameter is adjusted by a polynomial in
  \(s_1\) so that it maps to \(0\).  The conormal exact sequence identifies
  an integral differential modulo the jet with the differential of an element
  of the kernel, which is generated by \(s_1^{m+1}\) and \(s_2\).  Expanding
  the wedge of two such local expressions gives \(F\in(s_1^m,s_2)\).  A unit
  \(F\) forces \(m=0\).
\end{proof}

\begin{lemma}[Individual branch bound]
  \label{lem:cp_individual_branch_bound}
  \lean{MR4665779.IndividualBranchBound}
  \uses{lem:cp_branches_package, lem:cp_local_integral_expressions}
  % SOURCE: Caro--Pasten, Lemma 4.12 (LemmaIndividual) (read from references/caro-pasten-chabauty-coleman-surfaces-source/MR4665779 - A Chabauty-Coleman bound for surfaces.tex)
  % SOURCE QUOTE: $$\ord_y (\nu_j^\bullet(\omega))\ge \min\{m, \ord_{s_1}(\pi(F_{j,y})) -1\}.$$
  \textit{Source: Caro--Pasten, Lemma 4.12.}
  In the setup of the overdetermined bound, let \(F_{j,y}\) be the completed
  local equation of the branch of \(C_j\) corresponding to
  \(y\in\nu_j^{-1}(x)\).  If the jet \(\phi\) is \(\omega\)-integral, then
  \[
    \operatorname{ord}_y(\nu_j^\bullet\omega)
    \ge \min\{m,\operatorname{ord}_{s_1}(\pi(F_{j,y}))-1\}.
  \]
\end{lemma}
\begin{proof}
  Pull \(\omega\) back to the completed local ring of the normalized branch.
  The local expression for \(\omega\) from
  \(\ref{lem:cp_local_integral_expressions}\) gives terms involving
  \(h_1^m h_1'\), \(h_2h_1'\), and \(h_2'\), where
  \(h_i\) are the images of \(s_i\).  Płoski's estimate bounds
  \(\operatorname{ord}_t h_2\) below by
  \(\operatorname{ord}_{s_1}\pi(F_{j,y})\), and formal differentiation loses
  at most one order.  Taking the minimum of the resulting lower bounds gives
  the claim.
\end{proof}

\begin{theorem}[Overdetermined \(\omega\)-integrality bound]
  \label{thm:cp_overdetermined_bound}
  \lean{MR4665779.OverdeterminedBound}
  \uses{lem:cp_individual_branch_bound, lem:cp_local_integral_expressions}
  % SOURCE: Caro--Pasten, Theorem 4.4 (read from references/caro-pasten-chabauty-coleman-surfaces-source/MR4665779 - A Chabauty-Coleman bound for surfaces.tex)
  % SOURCE QUOTE: $$m\le \sum_{j=1}^\ell \sum_{y\in \nu_j^{-1}(x)} a_j\cdot \left(\ord_y (\nu_j^\bullet (\omega_0))+1 \right).$$
  \textit{Source: Caro--Pasten, Theorem 4.4.}
  Let \(S\) be a smooth irreducible surface over an algebraically closed field,
  \(x\in S(k)\), and \(\omega_1,\omega_2\) global \(1\)-forms with
  \(\omega_1\wedge\omega_2\ne0\).  Write
  \(\operatorname{div}(\omega_1\wedge\omega_2)=\sum a_jC_j\), and let
  \(\nu_j:\widetilde C_j\to S\) be normalization followed by inclusion.  If
  a length-\(m\) infinitesimal disk supported at \(x\) is integral for both
  \(\omega_1\) and \(\omega_2\), then for every
  \(\omega_0=c_1\omega_1+c_2\omega_2\),
  \[
    m\le \sum_j\sum_{y\in\nu_j^{-1}(x)}
      a_j\bigl(\operatorname{ord}_y(\nu_j^\bullet\omega_0)+1\bigr).
  \]
\end{theorem}
\begin{proof}
  If \(x\) is outside the divisor, the local expression lemma gives \(m=0\).
  Otherwise, apply the individual branch bound to \(\omega_0\).  If any branch
  has \(\operatorname{ord}_{s_1}\pi(F_{j,y})\ge m+1\), its summand already
  exceeds \(m\).  If all branch orders are at most \(m\), multiplying the
  branch equations shows that the right side is at least
  \(\operatorname{ord}_{s_1}\pi(F)\), and the local divisor expression gives
  \(\operatorname{ord}_{s_1}\pi(F)\ge m\).
\end{proof}

\section{Differentials on abelian varieties}

\begin{lemma}[Ample effective divisors on abelian varieties]
  \label{lem:cp_ample_abvar_divisor}
  \lean{MR4665779.AmpleAbelianVarietyDivisor}
  \uses{def:cp_geometry_interfaces}
  % SOURCE: Caro--Pasten, Lemma 5.1 (read from references/caro-pasten-chabauty-coleman-surfaces-source/MR4665779 - A Chabauty-Coleman bound for surfaces.tex)
  % SOURCE QUOTE: If $\supp(D)$ contains no translate of positive dimensional abelian subvarieties of $A$, then $D$ is ample.
  \textit{Source: Caro--Pasten, Lemma 5.1.}
  Let \(D\) be a nonzero effective divisor on an abelian variety \(A\).  If
  \(\operatorname{supp}(D)\) contains no translate of a positive-dimensional
  abelian subvariety, then \(D\) is ample.
\end{lemma}
\begin{proof}
  The stabilizer \(H(D)=\{x:\tau_x^*D=D\}\) is a closed subgroup of \(A\).
  For any point of the support, the translate of \(H(D)\) through that point
  lies in \(\operatorname{supp}(D)\).  The hypothesis therefore forces
  \(H(D)\) to be finite, and Mumford's ampleness criterion for divisors on
  abelian varieties applies.
\end{proof}

\begin{lemma}[Degree and trace package for curves in abelian varieties]
  \label{lem:cp_abvar_degree_trace_package}
  \lean{MR4665779.AbelianVarietyDegreeTracePackage}
  \uses{def:cp_geometry_interfaces}
  % SOURCE: Caro--Pasten, Section 5 (\label{SecRestrictions}) and Lemma 5.2 (LemmaSquare) (read from references/caro-pasten-chabauty-coleman-surfaces-source/MR4665779 - A Chabauty-Coleman bound for surfaces.tex)
  % SOURCE QUOTE: There is a monic polynomial $Q_{C,D}(t)\in \Z[t]$ of degree $n$ such that $P_{\alpha_A(C,D)}(t)=Q_{C,D}(t)^2$.
  \textit{Source: Caro--Pasten, Section 5 and Lemma 5.2.}
  For a curve \(C\) generating an abelian variety \(A\) of dimension \(n\) and
  an ample divisor \(D\), the endomorphism \(\alpha_A(C,D)\) has
  \(P_{\alpha_A(C,D)}(t)=Q_{C,D}(t)^2\) for a monic degree-\(n\) polynomial
  whose roots are real and positive.  Moreover,
  \(\operatorname{tr}(\alpha_A(C,D))=2\deg(C\cdot D)\), and
  \(\alpha_A(C,D)=\nu_{C/A,*}\nu_{C/A}^*\phi_D\).
\end{lemma}
\begin{proof}
  The trace formula and the factorization through the Jacobian of the
  normalization are Matsusaka--Morikawa identities.  Debarre's positivity
  theorem for the generated curve and ample divisor gives the square
  characteristic polynomial and positive real roots.
\end{proof}

\begin{lemma}[Degrees for surfaces in abelian threefolds]
  \label{lem:cp_degrees_surface_threefold}
  \lean{MR4665779.SurfaceThreefoldDegrees}
  \uses{def:cp_geometry_interfaces}
  % SOURCE: Caro--Pasten, Lemma 5.3 (read from references/caro-pasten-chabauty-coleman-surfaces-source/MR4665779 - A Chabauty-Coleman bound for surfaces.tex)
  % SOURCE QUOTE: $$\deg_X(X)=\cdeg_X(X)=\deg(X^3)_A=c_1^2(X).$$
  \textit{Source: Caro--Pasten, Lemma 5.3.}
  If \(X\) is a smooth projective surface contained in an abelian threefold
  \(A\), then \(X|_X\) is a canonical divisor of \(X\), and
  \[
    \deg_X(X)=\operatorname{cdeg}_X(X)=\deg(X^3)_A=c_1^2(X).
  \]
\end{lemma}
\begin{proof}
  The canonical sheaf of \(A\) is trivial, so adjunction identifies \(X|_X\)
  with a canonical divisor of \(X\).  The displayed equalities are then the
  same intersection number computed in \(A\) and on \(X\).
\end{proof}

\begin{lemma}[Injectivity of restricted one-forms]
  \label{lem:cp_injectivity_one_forms}
  \lean{MR4665779.InjectivityOneForms}
  \uses{lem:cp_abvar_degree_trace_package, lem:cp_degrees_surface_threefold, lem:cp_canonical_genus_bound}
  % SOURCE: Caro--Pasten, Lemmas 5.4--5.6 (read from references/caro-pasten-chabauty-coleman-surfaces-source/MR4665779 - A Chabauty-Coleman bound for surfaces.tex)
  % SOURCE QUOTE: Then the map $\nu_{C/A}^\bullet : H^0(A,\Omega^1_{A/k})\to H^0(\widetilde{C}, \Omega^1_{\widetilde{C}/k})$ is injective.
  \textit{Source: Caro--Pasten, Lemmas 5.4--5.6.}
  In positive characteristic, if a curve \(C\subset A\) generates \(A\) and
  \(p>\frac{n!}{n^n}\frac{\deg_H(C)^n}{\deg(H^n)}\), then
  \(H^0(A,\Omega^1_A)\to H^0(\widetilde C,\Omega^1_{\widetilde C})\) is
  injective.  For curves \(C\) in canonical divisors of the surface \(X\), the
  same argument gives either injectivity or a kernel of dimension at most
  \(1\) under the two Caro--Pasten hypotheses.  Under the surface hypotheses,
  restriction \(H^0(A,\Omega^1_A)\to H^0(X,\Omega^1_X)\) is injective.
\end{lemma}
\begin{proof}
  The degree--trace package gives an isogeny
  \(\beta:A^\vee\to A\) whose degree is a square \(M^2\), with
  \(M\) bounded by the arithmetic--geometric mean in terms of
  \(\deg_H(C)\).  The numerical hypothesis gives \(p\nmid M\), so \(\beta\)
  is separable and pullback of invariant forms is injective.  For curves in a
  canonical divisor, either the generated abelian subvariety is all of \(A\),
  or in the threefold case an abelian surface; the canonical genus bound
  controls the surface case and leaves a one-dimensional kernel.  For surfaces,
  intersect \(X\) with a very ample divisor to reduce to a smooth generating
  curve.
\end{proof}

\begin{lemma}[Nonvanishing of restricted two-forms]
  \label{lem:cp_nonvanishing_two_forms}
  \lean{MR4665779.NonvanishingTwoForms}
  \uses{lem:cp_injectivity_one_forms, lem:cp_degrees_surface_threefold, lem:cp_ample_abvar_divisor}
  % SOURCE: Caro--Pasten, Lemma 5.7 (read from references/caro-pasten-chabauty-coleman-surfaces-source/MR4665779 - A Chabauty-Coleman bound for surfaces.tex)
  % SOURCE QUOTE: Then $\iota^\bullet (\omega_1)\wedge \iota^\bullet(\omega_2)\in H^0(X,\Omega^2_{X/k})$ is not the zero section.
  \textit{Source: Caro--Pasten, Lemma 5.7.}
  Let \(X\subset A\) be a smooth surface in an abelian variety of dimension
  \(n\ge3\) over a positive-characteristic algebraically closed field.  If
  \(\omega_1,\omega_2\) are independent invariant differentials on \(A\), and
  either the threefold/no-elliptic-curve bound
  \(p>(128/9)c_1^2(X)^2\) or the geometrically simple ample-divisor bound
  holds, then the restricted \(2\)-form
  \(\iota^\bullet\omega_1\wedge\iota^\bullet\omega_2\) is not zero on \(X\).
\end{lemma}
\begin{proof}
  First the surface restriction of one-forms is injective, so the two
  restrictions are \(k\)-linearly independent.  If their wedge vanished, one
  would have \(u_1=f u_2\) for a nonconstant rational function.  A component
  \(C\) of the zero divisor of \(f\) is then integral for \(u_1\).  An
  auxiliary smooth curve \(D\sim3H|_X\) bounds \(H\cdot C\) by the degree of
  the nonzero pullback of \(u_1\) to \(D\).  The one-form injectivity theorem
  then contradicts \(\nu_{C/A}^\bullet(\omega_1)=0\), except in the abelian
  surface subcase of the threefold hypothesis; there the quotient by the
  generated abelian surface produces an elliptic curve and forces \(u_2\) to
  vanish on \(C\) as well, contradicting the one-dimensional kernel.
\end{proof}

\section{Zeros of \texorpdfstring{\(p\)}{p}-adic power series}

\begin{lemma}[One- and multi-variable zero estimates]
  \label{lem:cp_zero_estimates}
  \lean{MR4665779.ZeroEstimates}
  \uses{def:cp_notation}
  % SOURCE: Caro--Pasten, Lemmas 6.1 and 6.2 (read from references/caro-pasten-chabauty-coleman-surfaces-source/MR4665779 - A Chabauty-Coleman bound for surfaces.tex)
  % SOURCE QUOTE: If $h\ne 0$ and $\rho_h>0$, then for each $0<r<\rho_h$ we have $\nn_0(h,r)\le \nu(h,r)$.
  \textit{Source: Caro--Pasten, Lemmas 6.1 and 6.2.}
  For a nonzero one-variable \(p\)-adic power series \(h\), the number of
  zeros in a closed disk of radius \(r\) is at most the largest index attaining
  the Gauss norm at \(r\).  If a multivariable series \(H\) has coefficients
  bounded by \(M^{|\alpha|-1}\), and some homogeneous part of degree at most
  \(N\) evaluates to norm at least \(1\) on a unit vector \(u\), then along the
  line \(zu\),
  \[
    n_0(H_u,r)\le
    \left(N-\frac{\log M}{\log(r^{-1})}\right)
    \left(1-\frac{\log M}{\log(r^{-1})}\right)^{-1}
  \]
  for \(0<r<M^{-1}\).
\end{lemma}
\begin{proof}
  The first statement is the Newton polygon/Gauss norm zero bound.  For the
  second, the coefficient estimate implies convergence radius at least
  \(M^{-1}\).  If \(m>(N-\lambda)/(1-\lambda)\), with
  \(\lambda=\log M/\log(r^{-1})\), then every degree-\(m\) homogeneous term is
  strictly smaller than the nonzero contribution of degree at most \(N\).
  Thus the Gauss norm of \(H_u\) is attained by a degree no larger than
  \(\lfloor(N-\lambda)/(1-\lambda)\rfloor\), and the one-variable bound applies.
\end{proof}

\section{Exponential and logarithm}

\begin{lemma}[Local parameters and formal multiplication]
  \label{lem:cp_logexp_local_parameters}
  \lean{MR4665779.LogExpLocalParameters}
  \uses{def:cp_notation, def:cp_geometry_interfaces}
  % SOURCE: Caro--Pasten, Lemmas 7.1--7.4 (read from references/caro-pasten-chabauty-coleman-surfaces-source/MR4665779 - A Chabauty-Coleman bound for surfaces.tex)
  % SOURCE QUOTE: There are $t_j\in \mfrak_{\Acal,x_0}$ for $1\le j\le n$ such that $t_1,...,t_n$ are $R$-local parameters along $\Zcal$.
  \textit{Source: Caro--Pasten, Lemmas 7.1--7.4.}
  For an abelian scheme \(\mathcal A/R\), a section, and a smooth regular
  closed subscheme through it, there exist \(R\)-local parameters along the
  section whose first coordinates cut out the subscheme.  In such parameters,
  the formal multiplication-by-\(m\) expansions have coefficients in \(R\);
  multiplication by \(1\) has linear part \(t_j\), and the terms of degree
  \(<m\) in the homogeneous piece \(\Psi_j^{[m]}\) vanish.
\end{lemma}
\begin{proof}
  The parameter choice is the SGA smooth-local normal form around a section
  and a smooth closed subscheme.  Integrality of the formal group law over
  \(R\) gives integral multiplication series.  The assertions about the linear
  part and degree filtration are immediate from \([1]=\mathrm{id}\) and the
  standard formal-group identity for multiplication by \(m\).
\end{proof}

\begin{lemma}[Convergence and linearization of logarithm and exponential]
  \label{lem:cp_logexp_convergence_linearization}
  \lean{MR4665779.LogExpConvergenceLinearization}
  \uses{lem:cp_logexp_local_parameters}
  % SOURCE: Caro--Pasten, Lemmas 7.5--7.7 (read from references/caro-pasten-chabauty-coleman-surfaces-source/MR4665779 - A Chabauty-Coleman bound for surfaces.tex)
  % SOURCE QUOTE: The maps $\widetilde{\Log}^\ttt$ and $\widetilde{\Exp}^\ttt$ ... are isomorphisms of Lie groups, inverse to each other.
  \textit{Source: Caro--Pasten, Lemmas 7.5--7.7.}
  In local parameters \(\mathbf t\), the logarithm coefficients satisfy
  \(|b_{j,\alpha}|\le\|\alpha\|^{[K:\mathbb Q_p]}\), and the exponential
  coefficients satisfy
  \(m!c_{j,\alpha}\in R\) and
  \(|c_{j,\alpha}|\le p^{[K:\mathbb Q_p](m-1)/(p-1)}\) for
  \(m=\|\alpha\|\).  If
  \(p>\max\{\mathfrak e+1,\exp(\mathfrak e/\exp(1))\}\), the logarithm and
  exponential converge on the residue kernel and identify it, as a Lie group,
  with the additive ball \(\mathfrak p^n\).
\end{lemma}
\begin{proof}
  The logarithm coefficient bound comes from the integral multiplication
  series and denominators \(m^{-1}\).  The exponential bound follows from
  \(m!c_{j,\alpha}\in R\) and Legendre's estimate for \(|m!|\).  The two
  numerical hypotheses ensure both series converge on the radius \(q^{-1}\)
  ball and map it to itself.  The formal-group identities
  \(\log\circ\exp=\mathrm{id}\) and \(\exp\circ\log=\mathrm{id}\) then give
  inverse analytic Lie-group isomorphisms.
\end{proof}

\section{Analytic \texorpdfstring{\(1\)}{1}-parameter subgroups}

\begin{definition}[Analytic one-parameter subgroup]
  \label{def:cp_one_parameter_subgroup}
  \lean{MR4665779.OneParameterSubgroup}
  \uses{lem:cp_logexp_convergence_linearization}
  % SOURCE: Caro--Pasten, Section 8 (\label{Sec1PS}, subsection ``Definitions''; the definition is unnumbered prose) (read from references/caro-pasten-chabauty-coleman-surfaces-source/MR4665779 - A Chabauty-Coleman bound for surfaces.tex)
  % SOURCE QUOTE: A \emph{$1$-parameter subgroup} of $\Acal$ is a pair $\gamma=(\ttt,\uu)$ where $\ttt=(t_1,...,t_n)$ is a choice of $R$-local parameters at $e\in A(K)$ and $\uu\in K^n$ satisfies $|\uu|=1$.
  \textit{Source: Caro--Pasten, Section 8 (Sec1PS), ``Definitions''.}
  A one-parameter subgroup of \(\mathcal A\) is a pair
  \(\gamma=(\mathbf t,u)\), where \(\mathbf t\) is a choice of \(R\)-local
  parameters at the identity and \(u\in K^n\) has norm \(1\).  It defines
  \[
    \widetilde\gamma(z)=\widetilde{\operatorname{Exp}}^{\mathbf t}(zu)
      :\mathfrak p\to U_e
  \]
  and a hyperplane
  \(\mathcal H(\gamma)=\ker(\epsilon^u\circ
  (\widetilde{\operatorname{Exp}}^{\mathbf t})^\bullet)\) in
  \(H^0(A,\Omega^1_{A/K})\).
\end{definition}

\begin{lemma}[Hyperplanes and images of one-parameter subgroups]
  \label{lem:cp_one_parameter_hyperplanes}
  \lean{MR4665779.OneParameterHyperplanes}
  \uses{def:cp_one_parameter_subgroup}
  % SOURCE: Caro--Pasten, Lemmas 8.1--8.6 (read from references/caro-pasten-chabauty-coleman-surfaces-source/MR4665779 - A Chabauty-Coleman bound for surfaces.tex)
  % SOURCE QUOTE: The rule $\gamma\mapsto \Hcal(\gamma)$ defines a bijection between equivalence classes of $1$-parameter subgroups of $\Acal$ and $K$-linear hyperplanes.
  \textit{Source: Caro--Pasten, Lemmas 8.1--8.6.}
  The differential pullback of \(\widetilde\gamma\) has kernel
  \(\mathcal H(\gamma)\).  The image of \(\gamma\) is a closed saturated
  rank-one \(R\)-submodule of the residue kernel under logarithm.  Two
  one-parameter subgroups are equivalent exactly when their images agree,
  exactly when their images meet nontrivially, and equivalence classes are in
  bijection with \(K\)-linear hyperplanes of \(H^0(A,\Omega^1_{A/K})\).
  Every nonidentity element of \(U_e\) lies in the image of a unique such class.
\end{lemma}
\begin{proof}
  Under logarithm, the image of \(\gamma\) is \(\mathfrak p u\), so equality
  and nontrivial intersection of images are equivalent to scaling \(u\) by a
  unit after changing coordinates by the induced element of \(\operatorname{GL}_n(R)\).
  The hyperplane is the kernel of the linear functional evaluating invariant
  differentials in the direction \(u\).  Every hyperplane is the kernel of such
  a functional after choosing coordinates, and every nonzero vector in
  \(\mathfrak p^n\) normalizes to a unit vector \(u\).
\end{proof}

\begin{lemma}[Infinitesimal one-parameter subgroups and reduction]
  \label{lem:cp_infinitesimal_one_parameter}
  \lean{MR4665779.InfinitesimalOneParameterSubgroup}
  \uses{def:cp_one_parameter_subgroup, lem:cp_logexp_convergence_linearization}
  % SOURCE: Caro--Pasten, Lemmas 8.7 and 8.8 (read from references/caro-pasten-chabauty-coleman-surfaces-source/MR4665779 - A Chabauty-Coleman bound for surfaces.tex)
  % SOURCE QUOTE: If $m<p$, the morphism $\tilde{\gamma}_m:V_m\to A$ extends to an $R$-morphism $\Vcal_m\to \Acal$.
  \textit{Source: Caro--Pasten, Lemmas 8.7 and 8.8.}
  The \(m\)-th truncation \(\widetilde\gamma_m:\operatorname{Spec}K[z]/(z^{m+1})\to A\)
  is a closed immersion.  If \(m<p\), it extends to
  \(\mathcal V_m=\operatorname{Spec}R[z]/(z^{m+1})\to\mathcal A\), and its
  special fibre
  \(V'_m=\operatorname{Spec}\mathbf k[z]/(z^{m+1})\to A'\) is also a closed
  immersion.
\end{lemma}
\begin{proof}
  For \(m<p\), the exponential coefficient bounds make the truncated power
  series integral over \(R\), giving the integral model and special fibre.
  The linear term of the exponential is \(u_jz\) in some coordinate with
  \(u_j\in R^\times\), so the induced map on local rings contains an element
  \(u_jz+O(z^2)\), which generates the truncated polynomial algebra.  Hence
  the maps are closed immersions.
\end{proof}

\begin{lemma}[One-parameter subgroups are \(\omega\)-integral along their hyperplane]
  \label{lem:cp_one_parameter_omega_integral}
  \lean{MR4665779.OneParameterOmegaIntegral}
  \uses{def:cp_omega_integral, lem:cp_omega_integral_functorial_local, lem:cp_infinitesimal_one_parameter, lem:cp_one_parameter_hyperplanes}
  % SOURCE: Caro--Pasten, Lemmas 8.9--8.11 (read from references/caro-pasten-chabauty-coleman-surfaces-source/MR4665779 - A Chabauty-Coleman bound for surfaces.tex)
  % SOURCE QUOTE: Then $\tilde{\gamma}_m:V_m\to A$ is $\omega$-integral.
  \textit{Source: Caro--Pasten, Lemmas 8.9--8.11.}
  If \(\omega\in\mathcal H(\gamma)\), then every infinitesimal truncation
  \(\widetilde\gamma_m\) is \(\omega\)-integral.  If \(m\le p-2\) and
  \(\omega\) extends to an integral differential on \(\mathcal A\), then the
  special-fibre truncation \(\widetilde\gamma'_m\) is integral for the reduced
  differential.
\end{lemma}
\begin{proof}
  The analytic pullback of \(\omega\) by \(\widetilde\gamma\) vanishes because
  \(\omega\in\mathcal H(\gamma)\).  The compatibility lemma compares this
  analytic pullback with the algebraic pullback by each truncation.  For
  \(m\le p-2\), the module of differentials on \(R[z]/(z^{m+1})\) is free over
  \(R\), so vanishing after base change to \(K\) descends to \(R\) and then
  reduces modulo \(\mathfrak p\).
\end{proof}

\section{The main result}

\begin{lemma}[Rank-one reduction]
  \label{lem:cp_rank_one_reduction}
  \lean{MR4665779.RankOneReduction}
  \uses{def:cp_notation, def:cp_geometry_interfaces}
  % SOURCE: Caro--Pasten, Lemma 9.2 (read from references/caro-pasten-chabauty-coleman-surfaces-source/MR4665779 - A Chabauty-Coleman bound for surfaces.tex)
  % SOURCE QUOTE: It suffices to prove Theorem \ref{ThmMain} under the assumption $\rk(G)=1$.
  \textit{Source: Caro--Pasten, Lemma 9.2.}
  In Theorem \(\ref{thm:cp_main_local}\), it suffices to treat the case
  \(\operatorname{rank}G=1\).
\end{lemma}
\begin{proof}
  If \(G\) has rank \(0\), enlarge it by a non-torsion point of the uncountable
  group \(A(K)\).  The resulting group has rank \(1\), and its closure contains
  the original closure, so the rank-one bound implies the rank-zero bound.
\end{proof}

\begin{lemma}[Consequences of the local hypotheses]
  \label{lem:cp_main_hypotheses}
  \lean{MR4665779.MainHypotheses}
  \uses{lem:cp_general_type_package, lem:cp_ample_abvar_divisor}
  % SOURCE: Caro--Pasten, Lemma 9.3 (read from references/caro-pasten-chabauty-coleman-surfaces-source/MR4665779 - A Chabauty-Coleman bound for surfaces.tex)
  % SOURCE QUOTE: $X$ and $X'$ are surfaces of general type and $1\le c_1^2(X')=c_1^2(X)<(p-2)/3$.
  \textit{Source: Caro--Pasten, Lemma 9.3.}
  Under either local hypothesis of the main theorem, \(X\) and \(X'\) are
  surfaces of general type and
  \(1\le c_1^2(X')=c_1^2(X)<(p-2)/3\).  In the threefold case \(X'\) is ample
  on \(A'\); in the geometrically simple case the chosen ample divisor
  specializes with the same intersection and canonical degrees.
\end{lemma}
\begin{proof}
  Simplicity or absence of elliptic curves rules out the abelian-surface and
  elliptic-curve exceptions, so the general-type criterion applies.  Good
  reduction and the Katsura--Ueno specialization results preserve general type
  and \(c_1^2\).  The numerical hypotheses imply the displayed upper bound.
  In the threefold case, the effective divisor \(X'\) has support containing no
  translate of a positive-dimensional abelian subvariety, hence is ample.  In
  the simple case, specialization of numerical equivalence and intersection
  products preserves the relevant degrees.
\end{proof}

\begin{lemma}[Genus bound for curves in the special fibre]
  \label{lem:cp_special_fibre_curve_genus}
  \lean{MR4665779.SpecialFibreCurveGenus}
  \uses{lem:cp_main_hypotheses}
  % SOURCE: Caro--Pasten, Lemma 9.5 (read from references/caro-pasten-chabauty-coleman-surfaces-source/MR4665779 - A Chabauty-Coleman bound for surfaces.tex)
  % SOURCE QUOTE: Let $C$ be an irreducible curve in $X'_k$ defined over $k$. Then $\gfrak_g(C)\ge 2$.
  \textit{Source: Caro--Pasten, Lemma 9.5.}
  Every irreducible curve \(C\subset X'_k\) has geometric genus at least \(2\).
\end{lemma}
\begin{proof}
  A curve in an abelian variety has genus at least \(1\).  If the genus were
  \(1\), the normalization would be an elliptic curve and its nonconstant map
  to \(A'_k\) would have elliptic image contained in \(X'_k\), contradicting
  the hypotheses.
\end{proof}

\begin{lemma}[Choice of the Chabauty one-parameter subgroup]
  \label{lem:cp_choice_one_parameter_subgroup}
  \lean{MR4665779.ChoiceOneParameterSubgroup}
  \uses{lem:cp_rank_one_reduction, lem:cp_one_parameter_hyperplanes}
  % SOURCE: Caro--Pasten, Lemma 9.6 (read from references/caro-pasten-chabauty-coleman-surfaces-source/MR4665779 - A Chabauty-Coleman bound for surfaces.tex)
  % SOURCE QUOTE: There is a $1$-parameter subgroup $\gamma=(\ttt,\uu)$ for $A$ such that $\Gamma\cap U_e\subseteq \im(\gamma)$.
  \textit{Source: Caro--Pasten, Lemma 9.6.}
  For rank-one \(G\), and for any local parameters at the identity, there is a
  one-parameter subgroup \(\gamma\) such that
  \(\Gamma\cap U_e\subseteq\operatorname{im}\gamma\).  Its equivalence class is
  uniquely determined by the condition that \(\operatorname{im}\gamma\cap\Gamma\)
  strictly contains the identity.
\end{lemma}
\begin{proof}
  Choose a nonidentity point of \(\Gamma\cap U_e\).  The one-parameter
  hyperplane package puts it in a unique one-parameter subgroup up to
  equivalence.  Since \(\Gamma\cap U_e\) is the closure of a rank-one subgroup
  and \(\operatorname{im}\gamma\) is saturated and closed under logarithm, the
  nontrivial intersection forces the whole \(\Gamma\cap U_e\) into
  \(\operatorname{im}\gamma\).
\end{proof}

\begin{lemma}[Choice of two integral differentials]
  \label{lem:cp_choice_differentials}
  \lean{MR4665779.ChoiceDifferentials}
  \uses{lem:cp_choice_one_parameter_subgroup}
  % SOURCE: Caro--Pasten, Lemma 9.7 (read from references/caro-pasten-chabauty-coleman-surfaces-source/MR4665779 - A Chabauty-Coleman bound for surfaces.tex)
  % SOURCE QUOTE: There are $\omega_1,\omega_2\in \Hcal$ such that ... The differentials $\omega'_1,\omega'_2$ ... are $\kk$-linearly independent.
  \textit{Source: Caro--Pasten, Lemma 9.7.}
  Let \(\mathcal H\subset H^0(A,\Omega^1_{A/K})\) be the hyperplane associated
  to \(\Gamma\).  There exist \(K\)-linearly independent
  \(\omega_1,\omega_2\in\mathcal H\) extending to
  \(H^0(\mathcal A,\Omega^1_{\mathcal A/R})\) whose reductions
  \(\omega'_1,\omega'_2\) are \(\mathbf k\)-linearly independent.
\end{lemma}
\begin{proof}
  The \(R\)-module \(H^0(\mathcal A,\Omega^1_{\mathcal A/R})\) is free of rank
  \(n\).  The intersection with a \(K\)-hyperplane has rank \(n-1\ge2\).
  Choose two primitive independent vectors in this lattice; primitivity and
  independence prevent their reductions from becoming dependent over
  \(\mathbf k\).
\end{proof}

\begin{lemma}[The special-fibre canonical divisor from differentials]
  \label{lem:cp_special_canonical_divisor}
  \lean{MR4665779.SpecialCanonicalDivisor}
  \uses{lem:cp_choice_differentials, lem:cp_nonvanishing_two_forms, lem:cp_main_hypotheses, lem:cp_special_fibre_curve_genus, lem:cp_general_type_package}
  % SOURCE: Caro--Pasten, Lemmas 9.8--9.9 (read from references/caro-pasten-chabauty-coleman-surfaces-source/MR4665779 - A Chabauty-Coleman bound for surfaces.tex)
  % SOURCE QUOTE: The $2$-form $u_1\wedge u_2$ ... is not the zero section.
  \textit{Source: Caro--Pasten, Lemmas 9.8--9.9.}
  Let \(u_i\) be the restrictions of \(\omega'_i\) to \(X'\).  Then
  \(u_1\wedge u_2\) is nonzero.  Its divisor \(D\) on \(X'\) is effective,
  ample, defined over \(\mathbf k\), and canonical; in particular \(D>0\) and
  is nef.
\end{lemma}
\begin{proof}
  Nonvanishing is exactly the restricted two-form theorem applied to
  \(A'\) and \(X'\).  Since \(u_1\wedge u_2\) is a nonzero section of the
  canonical sheaf, its divisor is effective and canonical.  The main
  hypotheses give general type and exclude genus \(0\) and \(1\) curves in
  the support, so the general-type ampleness criterion makes \(D\) ample.
\end{proof}

\begin{lemma}[Choice of a nonintegral linear combination]
  \label{lem:cp_choice_w}
  \lean{MR4665779.ChoiceW}
  \uses{lem:cp_special_canonical_divisor, lem:cp_injectivity_one_forms, lem:cp_canonical_genus_bound}
  % SOURCE: Caro--Pasten, Lemma 9.10 (read from references/caro-pasten-chabauty-coleman-surfaces-source/MR4665779 - A Chabauty-Coleman bound for surfaces.tex)
  % SOURCE QUOTE: There is $b\in \kk$ such that the differential $w=u_1+bu_2$ ... satisfies that no map $\nu_j$ is $w$-integral.
  \textit{Source: Caro--Pasten, Lemma 9.10.}
  Writing \(D=\sum a_jC_j\) over \(k=\mathbf k^{alg}\), there exists
  \(b\in\mathbf k\) such that \(w=u_1+bu_2\) has nonzero pullback to every
  normalization \(\nu_j:\widetilde C_j\to X'_k\).
\end{lemma}
\begin{proof}
  Ampleness and the canonical genus bound imply that the number of components
  \(\ell\) is \(<p\).  If every \(b\in\mathbb F_p\) made some component
  \(w_b\)-integral, two distinct \(b,b'\) would vanish on the same component.
  Then the restrictions of both \(\omega'_1\) and \(\omega'_2\) would lie in
  the kernel for that curve, making the kernel at least two-dimensional,
  contradicting the injectivity or semi-injectivity theorem for one-forms.
\end{proof}

\begin{lemma}[Bounds for infinitesimal integral disks]
  \label{lem:cp_mx_bounds}
  \lean{MR4665779.MxBounds}
  \uses{thm:cp_overdetermined_bound, lem:cp_choice_w, lem:cp_special_fibre_curve_genus, lem:cp_branches_package, lem:cp_canonical_genus_bound, lem:cp_main_hypotheses}
  % SOURCE: Caro--Pasten, Lemmas 9.11--9.13 (read from references/caro-pasten-chabauty-coleman-surfaces-source/MR4665779 - A Chabauty-Coleman bound for surfaces.tex)
  % SOURCE QUOTE: For every $x\in X'(k)$ we have $m(x)\le p-3$.
  \textit{Source: Caro--Pasten, Lemmas 9.11--9.13.}
  Let \(m(x)\) be the supremum length of closed infinitesimal disks in
  \(X'_k\) supported at \(x\) and integral for both \(u_1,u_2\).  If
  \(x\notin\operatorname{supp}D\), then \(m(x)=0\).  If \(x\in\operatorname{supp}D\),
  then
  \[
    m(x)\le \sum_j\sum_{y\in\nu_j^{-1}(x)}
      a_j(\operatorname{ord}_y\nu_j^\bullet w+1),
  \]
  and for every \(x\), \(m(x)\le p-3\).
\end{lemma}
\begin{proof}
  The first two assertions are the overdetermined bound, using
  \(\omega_0=w\).  For the uniform \(p-3\) bound, sum the orders of the
  nonzero differentials \(\nu_j^\bullet w\), whose degrees are
  \(2g_g(C_j)-2\), and use the normalization-fibre bound
  \(\#\nu_j^{-1}(x)\le\delta(C_j,x)+1\le g_a(C_j)-1\).  The genus bound
  \(g_g(C_j)\ge2\), the canonical genus bound, and
  \(3c_1^2(X')<p-2\) give \(m(x)<p-2\).
\end{proof}

\begin{proposition}[Residue disk bound]
  \label{prop:cp_residue_disk_bound}
  \lean{MR4665779.ResidueDiskBound}
  \uses{lem:cp_mx_bounds, lem:cp_choice_one_parameter_subgroup, lem:cp_logexp_local_parameters, lem:cp_infinitesimal_one_parameter, lem:cp_one_parameter_omega_integral, lem:cp_zero_estimates}
  % SOURCE: Caro--Pasten, Proposition 9.14 and Lemmas 9.15--9.16 (read from references/caro-pasten-chabauty-coleman-surfaces-source/MR4665779 - A Chabauty-Coleman bound for surfaces.tex)
  % SOURCE QUOTE: $$\# \Gamma\cap X(K)\cap U_x\le 1+ m(x)\cdot  \left(1-\frac{\efrak}{p-1}\right)^{-1}.$$
  \textit{Source: Caro--Pasten, Proposition 9.14 and Lemmas 9.15--9.16.}
  For \(\xi\in\Gamma\cap X(K)\) with reduction \(x\in X'(\mathbf k)\),
  \[
    \#(\Gamma\cap X(K)\cap U_x)
    \le 1+m(x)\left(1-\frac{\mathfrak e}{p-1}\right)^{-1}.
  \]
  In particular, if \(m(x)=0\), the residue disk contains at most one point of
  \(\Gamma\cap X(K)\).
\end{proposition}
\begin{proof}
  Translate the situation so that \(\xi=e\); translation preserves \(\Gamma\),
  residue disks, invariant differentials, and \(m(x)\).  Choose local
  parameters whose first \(n-2\) coordinates cut out \(\mathcal X\), and choose
  the rank-one one-parameter subgroup whose image contains \(\Gamma\cap U_e\).
  Then points of \(\Gamma\cap X(K)\cap U_e\) are zeros of the one-variable
  series \(\operatorname{Exp}^{\mathbf t}_j(zu)\) for each defining coordinate.
  If all homogeneous coefficients of degrees \(\le m(e)\) had norm \(<1\),
  the reduced truncation would factor through \(X'\) and be integral for
  \(u_1,u_2\), contradicting the definition of \(m(e)\); hence some
  \(P_{j,N}^{\mathbf t}(u)\) has norm at least \(1\) with \(N\le m(e)+1\).
  The multivariable zero estimate with
  \(M=p^{[K:\mathbb Q_p]/(p-1)}\) and \(r=q^{-1}\) gives the displayed bound.
\end{proof}

\begin{theorem}[Local Chabauty--Coleman bound for surfaces]
  \label{thm:cp_main_local}
  \lean{MR4665779.MainLocal}
  \uses{lem:cp_rank_one_reduction, prop:cp_residue_disk_bound, lem:cp_mx_bounds, lem:cp_weil_bound_singular_curve, lem:cp_canonical_genus_bound, lem:cp_special_canonical_divisor}
  % SOURCE: Caro--Pasten, Theorem 9.1 (read from references/caro-pasten-chabauty-coleman-surfaces-source/MR4665779 - A Chabauty-Coleman bound for surfaces.tex)
  % SOURCE QUOTE: $$\# \Gamma\cap X(K) \le \# X'(\kk) + \left(1-\frac{\efrak}{p-1}\right)^{-1}\cdot (q+4q^{1/2}+3)\cdot c_1^2(X).$$
  \textit{Source: Caro--Pasten, Theorem 9.1.}
  Let \(\mathcal A/R\) be an abelian scheme of relative dimension \(n\ge3\),
  let \(\mathcal X\subset\mathcal A\) be smooth proper of relative dimension
  \(2\) with geometrically irreducible fibres, and let
  \(\Gamma\) be the \(p\)-adic closure of a finitely generated subgroup
  \(G\le A(K)\) of rank at most \(1\).  Assume
  \(p>\max\{\mathfrak e+1,\exp(\mathfrak e/\exp(1))\}\) and that either
  \begin{itemize}
    \item[(i)] \(n=3\), \(X\) is of general type, \(X'\) contains no elliptic
      curves over \(\mathbf k^{alg}\), and \(p>(128/9)c_1^2(X)^2\); or
    \item[(ii)] \(A'\) is simple over \(\mathbf k^{alg}\) and there is an ample
      divisor \(H\) on \(A\) such that
      \[
        p>\max\left\{3c_1^2(X)+2,\,
          \frac{n!\,(3\deg_H(X)+\operatorname{cdeg}_H(X))^n}
               {n^n\,\deg(H^n)}\right\}.
      \]
  \end{itemize}
  Then
  \[
    \#(\Gamma\cap X(K))\le
    \#X'(\mathbf k)+
    \left(1-\frac{\mathfrak e}{p-1}\right)^{-1}
    (q+4q^{1/2}+3)c_1^2(X).
  \]
\end{theorem}
\begin{proof}
  By rank-one reduction assume \(\operatorname{rank}G=1\).  Sum the residue
  disk bound over \(x\in X'(\mathbf k)\).  For \(x\notin D\), \(m(x)=0\), so
  these disks contribute exactly the main term.  For \(x\in D\), the
  overdetermined bound gives a sum of orders of \(\nu_j^\bullet w\) plus
  weighted normalization-fibre counts.  The order sum is bounded by
  \(2c_1^2(X)\) via the canonical genus bound.  The weighted fibre count is
  bounded by the modified Weil estimate and then again by the canonical genus
  bound, giving \((q+4q^{1/2}+1)c_1^2(X)\).  Together these contributions
  yield \((q+4q^{1/2}+3)c_1^2(X)\), multiplied by the residue-disk factor.
\end{proof}

\begin{theorem}[The \(\mathbb Q_p\) specialization]
  \label{thm:cp_main_qp}
  \lean{MR4665779.MainQp}
  \uses{thm:cp_main_local, lem:cp_main_hypotheses}
  % SOURCE: Caro--Pasten, Theorem 1.1 and Remark 9.4 (read from references/caro-pasten-chabauty-coleman-surfaces-source/MR4665779 - A Chabauty-Coleman bound for surfaces.tex)
  % SOURCE QUOTE: Theorem \ref{ThmMainIntro} follows from Theorem \ref{ThmMain}.
  \textit{Source: Caro--Pasten, Theorem 1.1 and Remark 9.4.}
  Over \(K=\mathbb Q_p\), with \(X\subset A\) of dimension \(n\ge3\) of good
  reduction and \(G\le A(\mathbb Q_p)\) of rank at most \(1\), assume that
  either
  \begin{itemize}
    \item[(i)] \(n=3\), \(X\) is of general type, \(X'\) contains no elliptic
      curves over \(\mathbb F_p^{alg}\), and \(p>(128/9)c_1^2(X)^2\); or
    \item[(ii)] \(A'\) is simple over \(\mathbb F_p^{alg}\),
      \(p>3c_1^2(X)+2\), and there is an ample divisor \(H\) on \(A\) such that
      \[
        p>\frac{n!\,(3\deg(H^2.X)+\deg(H.K_X))^n}{n^n\,\deg(H^n)}.
      \]
  \end{itemize}
  Then
  \[
    \#(X(\mathbb Q_p)\cap\Gamma)
    \le \#X'(\mathbb F_p)+
    \frac{p-1}{p-2}\bigl(p+4p^{1/2}+3\bigr)c_1^2(X).
  \]
\end{theorem}
\begin{proof}
  In the local theorem, \(\mathfrak e=1\), \(q=p\), and the preliminary
  numerical condition becomes \(p\ge3\).  The main hypotheses imply
  \(p\ge7\) because \(c_1^2(X)\ge1\), so no extra condition remains.  The
  displayed formula is the specialization of the local bound.
\end{proof}

\section{Applications for rational points}

\subsection{Rational points on surfaces}

\begin{theorem}[Number-field rational-point bound]
  \label{thm:cp_number_field_bound}
  \lean{MR4665779.NumberFieldBound}
  \uses{thm:cp_main_local}
  % SOURCE: Caro--Pasten, Theorem 10.1 (read from references/caro-pasten-chabauty-coleman-surfaces-source/MR4665779 - A Chabauty-Coleman bound for surfaces.tex)
  % SOURCE QUOTE: Then $X(F)$ is finite and we have $$\# X(F) \le \# X'(\kk) + \left(1-\frac{\efrak}{p-1}\right)^{-1}\cdot (q+4q^{1/2}+3)\cdot c_1^2(X).$$
  \textit{Source: Caro--Pasten, Theorem 10.1.}
  Let \(F\) be a number field and \(\mathfrak p\mid p\) a prime satisfying the
  local numerical condition.  Let \(X\subset A\) over \(F\) have good reduction
  at \(\mathfrak p\), with \(\dim A\ge3\), \(\operatorname{rank}A(F)\le1\),
  and one of the two Caro--Pasten geometric hypotheses after reduction.  Then
  \(X(F)\) is finite and satisfies the same bound as the local theorem with
  \(K=F_{\mathfrak p}\).
\end{theorem}
\begin{proof}
  Embed \(X(F)\) into \(X(F_{\mathfrak p})\cap\Gamma\), where \(\Gamma\) is
  the \(p\)-adic closure of \(A(F)\) in \(A(F_{\mathfrak p})\).  The
  Mordell--Weil rank condition gives \(\operatorname{rank}A(F)\le1\), so the
  local theorem applies after base change to the completion.
\end{proof}

\begin{theorem}[Rational surface bounds over \(\mathbb Q\)]
  \label{thm:cp_intro_surface_bounds}
  \lean{MR4665779.IntroSurfaceBounds}
  \uses{thm:cp_main_qp, thm:cp_number_field_bound, lem:cp_general_type_package}
  % SOURCE: Caro--Pasten, Theorems 1.5 and 1.7 (read from references/caro-pasten-chabauty-coleman-surfaces-source/MR4665779 - A Chabauty-Coleman bound for surfaces.tex)
  % SOURCE QUOTE: The following two results on rational points of surfaces are deduced from Theorem 1.1 by base change to Qp.
  \textit{Source: Caro--Pasten, Theorems 1.5 and 1.7.}
  The two stated \(\mathbb Q\)-surface theorems follow: for a general-type
  surface \(X\) in an abelian threefold with \(\operatorname{rank}A(\mathbb Q)\le1\)
  and no elliptic curves over \(\mathbb F_p^{alg}\) in the good reduction, where
  \(p>(128/9)c_1^2(X)^2\); and for a surface \(X\subset A\) of dimension
  \(n\ge3\) with \(\operatorname{rank}A(\mathbb Q)\le1\) and geometrically
  simple reduction \(A'\), where for an ample divisor \(H\) on \(A\),
  \[
    p>\max\left\{3c_1^2(X)+2,\,
      \frac{n!\,(3\deg(H^2.X)+\deg(H.K_X))^n}{n^n\,\deg(H^n)}\right\},
  \]
  the point set \(X(\mathbb Q)\) is finite and satisfies the explicit
  \(\mathbb F_p\)-point plus
  \(\frac{p-1}{p-2}(p+4p^{1/2}+3)c_1^2(X)\) bound.
\end{theorem}
\begin{proof}
  Both are the number-field theorem over \(F=\mathbb Q\).  In the threefold
  statement the hypotheses are exactly local hypothesis (i).  In the
  higher-dimensional statement they are exactly local hypothesis (ii).  The
  inclusion \(X(\mathbb Q)\subset X(\mathbb Q_p)\cap\Gamma\) transfers the
  local estimate to rational points.
\end{proof}

\subsection{Primes of geometrically simple reduction}

\begin{lemma}[Chavdarov density of geometrically simple reduction]
  \label{lem:cp_chavdarov_simple_reduction}
  \lean{MR4665779.ChavdarovSimpleReduction}
  \uses{def:cp_geometry_interfaces}
  % SOURCE: Caro--Pasten, Lemma 10.2 (read from references/caro-pasten-chabauty-coleman-surfaces-source/MR4665779 - A Chabauty-Coleman bound for surfaces.tex)
  % SOURCE QUOTE: There is a set of primes $\Pcal$ of density $1$ in the primes such that for every $p\in \Pcal$ the reduction of $A$ modulo $p$ is geometrically simple.
  \textit{Source: Caro--Pasten, Lemma 10.2.}
  If \(A/\mathbb Q\) is an abelian variety of odd dimension \(n\) with
  \(\operatorname{End}(A_\mathbb C)=\mathbb Z\), then a density-one set of
  primes has geometrically simple reduction for \(A\).
\end{lemma}
\begin{proof}
  Chavdarov's theorem (N. Chavdarov, ``The generic irreducibility of the numerator
  of the zeta function in a family of curves with large monodromy,'' Duke Math.\ J.\
  87 (1997), 151--180) establishes this for abelian varieties whose
  \(\ell\)-adic monodromy group is the full symplectic group.  For an abelian
  variety \(A/\mathbb Q\) of odd dimension \(n\) with
  \(\operatorname{End}(A_\mathbb C)=\mathbb Z\), the Tate conjecture (proved by
  Faltings) and the classification of maximally symplectic monodromy groups
  imply \(G_{\ell,A}=\operatorname{Sp}_{2n}\) for all but finitely many \(\ell\).
  Chavdarov's result then gives that for a density-one set of primes, the
  characteristic polynomial of Frobenius on \(T_\ell A\) is irreducible over
  \(\mathbb Q\), which implies geometric simplicity of the reduction.
\end{proof}

\begin{corollary}[Density-one threefold applicability]
  \label{cor:cp_dim3_density_application}
  \lean{MR4665779.Dim3DensityApplication}
  \uses{lem:cp_chavdarov_simple_reduction, thm:cp_intro_surface_bounds}
  % SOURCE: Caro--Pasten, Corollary 1.11 (CoroDim3) and proof in Section \label{SecApp2} (``Primes of geometrically simple reduction'') (read from references/caro-pasten-chabauty-coleman-surfaces-source/MR4665779 - A Chabauty-Coleman bound for surfaces.tex)
  % SOURCE QUOTE: Let $A$ be an abelian threefold over $\Q$ with $\rk A(\Q)\le 1$ and $\End (A_\C) = \Z$.
  \textit{Source: Caro--Pasten, Corollary 1.11 and Section SecApp2.}
  For an abelian threefold \(A/\mathbb Q\) with
  \(\operatorname{rank}A(\mathbb Q)\le1\) and
  \(\operatorname{End}(A_\mathbb C)=\mathbb Z\), a density-one set of primes
  has the property that every smooth geometrically irreducible projective
  surface \(X\subset A\) with good reduction at such a prime and satisfying
  \(p>(128/9)c_1^2(X)^2\) obeys the explicit rational-point bound.
\end{corollary}
\begin{proof}
  Apply Chavdarov to obtain geometrically simple reductions after discarding
  finitely many bad primes.  A surface in a geometrically simple abelian
  threefold reduction contains no elliptic curve, so the threefold rational
  surface theorem applies.
\end{proof}

\subsection{Quadratic points in curves}

\begin{lemma}[Intersection theory of \(C^{(2)}\)]
  \label{lem:cp_symmetric_square_intersections}
  \lean{MR4665779.SymmetricSquareIntersections}
  \uses{def:cp_geometry_interfaces}
  % SOURCE: Caro--Pasten, Lemmas 10.3 (LemmaKSymm2), 10.4 (LemmaThetaC2), 10.6 (LemmaThetaK) and Corollary 10.5 (LemmaSymm2GT) (read from references/caro-pasten-chabauty-coleman-surfaces-source/MR4665779 - A Chabauty-Coleman bound for surfaces.tex)
  % SOURCE QUOTE: $$K_{C^{(2)}}\equiv (2g-2)V-D/2.$$
  \textit{Source: Caro--Pasten, Lemmas 10.3, 10.4, 10.6 and Corollary 10.5.}
  For a smooth curve \(C\) of genus \(g\ge2\) in characteristic not \(2\),
  \[
    K_{C^{(2)}}\equiv(2g-2)V-\frac12D,\quad
    V^2=1,\quad D\cdot V=2,\quad D^2=4-4g,
  \]
  and \(c_1^2(C^{(2)})=(4g-9)(g-1)\).  If \(\theta\) is the pullback of a
  theta divisor under the Abel--Jacobi map, then
  \(D\equiv2((g+1)V-\theta)\), so
  \(K_{C^{(2)}}\equiv\theta+(g-3)V\); hence for \(g\ge3\),
  \(C^{(2)}\) has ample canonical divisor and is of general type.  Also
  \(\theta\cdot K_{C^{(2)}}=2g(g-2)\).
\end{lemma}
\begin{proof}
  Adjunction for the quotient \(C\times C\to C^{(2)}\), ramified along the
  diagonal, gives the canonical class and basic intersections.  Kouvidakis's
  numerical equivalence formula gives the relation with \(\theta\).  Since
  \(\theta\) is ample and \(V\) is effective, \(K\equiv\theta+(g-3)V\) is ample
  for \(g\ge3\).  Substituting the numerical equivalences gives the
  \(c_1^2\) and \(\theta\cdot K\) formulas.
\end{proof}

\begin{lemma}[Embedding \(C^{(2)}\) in the Jacobian]
  \label{lem:cp_symmetric_square_jacobian_embedding}
  \lean{MR4665779.SymmetricSquareJacobianEmbedding}
  \uses{lem:cp_symmetric_square_intersections}
  % SOURCE: Caro--Pasten, Lemmas 10.7 (LemmanonHypEll) and 10.8 (LemmaTheta) (read from references/caro-pasten-chabauty-coleman-surfaces-source/MR4665779 - A Chabauty-Coleman bound for surfaces.tex)
  % SOURCE QUOTE: If $C$ is not hyperelliptic, then $j$ is injective, in which case $C^{(2)}\simeq W_2$.
  \textit{Source: Caro--Pasten, Lemmas 10.7 and 10.8.}
  If \(C\) is not hyperelliptic, the Abel--Jacobi map
  \(j:C^{(2)}\to J\) is injective and identifies \(C^{(2)}\) with its image
  \(W_2\).  If \(C\) has an effective degree-\(2\) divisor over a field \(F\),
  then \(J\) has an \(F\)-defined ample divisor \(H\) with
  \(H\equiv2\Theta\).
\end{lemma}
\begin{proof}
  Noninjectivity of the degree-\(2\) Abel--Jacobi map is exactly the existence
  of a \(g^1_2\), i.e. hyperellipticity.  For the divisor \(H\), write the
  degree-\(2\) divisor as \(x_1+x_2\), possibly a conjugate pair, and add the
  two theta translates obtained from the \((g-1)\)-fold sums of
  \(j(V_{x_i})\); the sum descends to \(F\) and is numerically \(2\Theta\).
\end{proof}

\begin{corollary}[Quadratic point bound for nonhyperelliptic curves]
  \label{cor:cp_quadratic_points}
  \lean{MR4665779.QuadraticPoints}
  \uses{thm:cp_intro_surface_bounds, lem:cp_symmetric_square_intersections, lem:cp_symmetric_square_jacobian_embedding}
  % SOURCE: Caro--Pasten, Corollary 1.14 (CoroQuadratic) and proof in Section \label{SecApp3} (``Quadratic points in curves'') (read from references/caro-pasten-chabauty-coleman-surfaces-source/MR4665779 - A Chabauty-Coleman bound for surfaces.tex)
  % SOURCE QUOTE: Then $C^{(2)}(\Q)$ is finite and ...
  \textit{Source: Caro--Pasten, Corollary 1.14 and Section SecApp3.}
  Let \(C/\mathbb Q\) be smooth, geometrically irreducible, nonhyperelliptic
  over \(\mathbb Q^{alg}\), of genus \(g\ge3\), and suppose
  \(\operatorname{rank}J(\mathbb Q)\le1\).  If
  \(p>(8g-10)^g\) is a good prime such that \(C'\) is nonhyperelliptic and
  \(J'\) is geometrically simple, then \(C^{(2)}(\mathbb Q)\) is finite and
  \[
    \#C^{(2)}(\mathbb Q)\le \#(C')^{(2)}(\mathbb F_p)
      +\frac{p-1}{p-2}(p+4p^{1/2}+3)(4g-9)(g-1).
  \]
\end{corollary}
\begin{proof}
  If \(C^{(2)}(\mathbb Q)\) is nonempty, choose the corresponding degree-\(2\)
  divisor to define \(j:C^{(2)}\simeq W_2\subset J\).  Take
  \(H\equiv2\Theta\).  Poincaré formulas give
  \(\deg(H^g)=2^g g!\), \(\deg(H^2\cdot W_2)=4g(g-1)\), and
  \(\deg(H\cdot K_{W_2})=4g(g-2)\).  Hence the ample-divisor threshold in the
  surface theorem is \((8g-10)^g\), and the \(c_1^2\) term is
  \((4g-9)(g-1)\).  Apply the geometrically simple rational surface theorem.
\end{proof}

\begin{corollary}[Genus-three quadratic point bound]
  \label{cor:cp_quadratic_points_genus_three}
  \lean{MR4665779.QuadraticPointsGenusThree}
  \uses{thm:cp_intro_surface_bounds, lem:cp_symmetric_square_intersections, lem:cp_symmetric_square_jacobian_embedding}
  % SOURCE: Caro--Pasten, Corollary 1.15 (CoroQuadratic3) and proof in Section \label{SecApp3} (``Quadratic points in curves'') (read from references/caro-pasten-chabauty-coleman-surfaces-source/MR4665779 - A Chabauty-Coleman bound for surfaces.tex)
  % SOURCE QUOTE: Thus, $(128/9)c_1^2(X)^2= 512$ and the least prime $p>512$ is $521$.
  \textit{Source: Caro--Pasten, Corollary 1.15 and Section SecApp3.}
  If \(C/\mathbb Q\) has genus \(3\), is nonhyperelliptic, has
  \(\operatorname{rank}J(\mathbb Q)\le1\), has good reduction at
  \(p\ge521\), and \((C')^{(2)}\) contains no elliptic curves, then
  \[
    \#C^{(2)}(\mathbb Q)\le \#(C')^{(2)}(\mathbb F_p)
      +6\frac{p-1}{p-2}(p+4p^{1/2}+3)
      <\#(C')^{(2)}(\mathbb F_p)+7.1p .
  \]
\end{corollary}
\begin{proof}
  For \(g=3\), the symmetric-square intersection formula gives
  \(c_1^2(C^{(2)})=6\) and the canonical divisor is ample.  Thus the
  threefold rational surface theorem applies to \(W_2\simeq C^{(2)}\).  The
  threshold is \((128/9)\cdot6^2=512\), so the least prime satisfying it is
  \(521\), and the displayed numerical estimate is the specialization of the
  theorem's bound.
\end{proof}
